\section{Fuentes de señal auxiliares}

Además de las cuatro fuentes de cámara, el sistema incorpora varias fuentes de señal
auxiliares que permiten gestionar el estado de la emisión en momentos en los que la
mezcla principal no debe ser visible para el espectador, así como una entrada de audio
independiente para música de fondo.

\subsection{Fuente break e intro}

\texttt{break} e \texttt{intro} son fuentes de vídeo pregrabadas que el núcleo trata
exactamente igual que una cámara adicional. Se conectan a los puertos 10004 (break) y
10005 (intro) mediante bucles FFmpeg:

\begin{itemize}
  \item \textbf{break}: vídeo en bucle empleado como pantalla de pausa animada durante
  los descansos del evento. El realizador puede seleccionarlo como fuente activa del
  mezclador o bien activarlo automáticamente a través del mecanismo de stream blanker.
  \item \textbf{intro}: vídeo introductorio que se reproduce al inicio del evento o
  entre sesiones. Se implementa como un bucle continuo pre-cargado, de forma que esté
  siempre disponible para el realizador sin necesidad de iniciar un proceso adicional.
\end{itemize}

\subsection{Entrada de audio}

El sistema incorpora una entrada de audio independiente en el puerto 18000, pensada
para inyectar música de fondo durante los estados de pausa o fuera de emisión. Esta
entrada acepta un flujo Matroska de sólo audio (PCM S16LE, 48~kHz, estéreo) y su
nivel de salida puede controlarse desde el protocolo de control.

\subsection{Stream Blanker: estados de emisión}

El stream blanker es el mecanismo que permite interrumpir la señal de programa hacia el
exterior sin interrumpir el flujo de trabajo interno del realizador. Se implementa en
\texttt{voctocore/lib/streamblanker.py} como un elemento GStreamer adicional interpuesto
entre la mezcla y el puerto de salida de streaming (puerto 15000).

El sistema define tres estados de emisión, controlables desde la barra de herramientas
de voctogui:

\begin{itemize}
  \item \textbf{LIVE:} se emite la mezcla activa producida por voctocore. El realizador
        controla en todo momento lo que ve el espectador.
  \item \textbf{PAUSE:} se emite el vídeo de pausa procedente del puerto 17000 (fuente
        \texttt{stream-blanker-pause}), acompañado de la música de fondo del puerto 18000.
        Se usa durante los descansos del evento para que los espectadores vean una pantalla
        animada mientras el equipo de producción prepara la siguiente sesión.
  \item \textbf{NOSTREAM:} se emite un vídeo de ``fuera de emisión'' procedente del
        puerto 17001 (fuente \texttt{stream-blanker-nostream}), indicando al espectador
        que la transmisión ha finalizado o aún no ha comenzado.
\end{itemize}

La transición entre estados se realiza modificando los pesos del compositor de
GStreamer en tiempo real, evitando cortes bruscos en el stream de salida hacia el
encoder. Desde la interfaz gráfica, el realizador dispone de tres botones (LIVE, PAUSE,
NOSTREAM) en la barra de herramientas (\texttt{voctogui/lib/toolbar/streamblank.py}).

% --- Figura sugerida ---
% \begin{figure}[H]
%   \centering
%   \includegraphics[width=0.6\textwidth]{figuras/fsm_stream_blanker.png}
%   \caption{Máquina de estados del stream blanker (LIVE / PAUSE / NOSTREAM).}
%   \label{fig:fsm_stream_blanker}
% \end{figure}
